\documentclass[12pt]{article}
\usepackage{geometry}
\geometry{margin=1.4in}
\usepackage{setspace}
\setstretch{1.4}
\usepackage{amsmath}

\title{Ecology as Displacement\\
Species, Ecosystems, and the Cost of Extraction}
\author{Diego Rincón}
\date{}

\begin{document}

\maketitle

\begin{abstract}
An ecosystem is a web of relationships in equilibrium. Ground state $S_0$: each species exists at sustainable population, energy flows cyclically, diversity is maintained. Actual state $S^*$: humans extract resources at rate faster than regeneration, species are eliminated, diversity collapses, energy flows one-way (from earth to waste). The displacement is catastrophic and accelerating. The cost of displacement $D(\xi)$ is paid by non-human species (extinction, suffering, ecosystem collapse) and by humans (loss of ecosystem services, climate instability, eventual starvation). The cost of return $C_{\text{return}}$ is immense: stop extraction, restore degraded land, accept lower material living standards, wait decades for ecosystems to recover. Most humans and their governments will not pay this cost. Collapse is the likely outcome.
\end{abstract}

\section{The Ecosystem: Ground State}

In $S_0$, energy enters as sunlight. Plants capture it. Herbivores eat plants. Carnivores eat herbivores. Decomposers break down all into soil. Nutrients cycle. Species populations are stable (within natural variation). The system is resilient: if one species declines, others fill the niche.

Humans were part of this system for 200,000 years. We extracted resources (hunted, gathered, farmed), but at rates the system could regenerate.

\section{Extraction: Displacement}

50 years ago, extraction accelerated. Humans now:
- Fish oceans faster than they regenerate (90\% of fisheries collapsed)
- Cut forests faster than they regrow (50\% of forests gone)
- Extract fossil fuels (sunlight captured 300 million years ago, burned in 200 years)
- Farm monocultures (biodiversity destroyed for efficiency)
- Pollute (plastic, chemicals, radiation accumulate)
- Change climate (CO2 traps heat, disrupts weather patterns)

The displacement is from $S_0$ (sustainable extraction) to $S^*$ (extractive collapse).

\section{The Cost of Displacement}

Paid by non-human species:
- 68\% of wildlife populations have declined since 1970
- 1 million species face extinction
- Pollinators are crashing (bees, butterflies)
- Coral reefs are dying
- Ocean is acidifying and warming

Paid by humans:
- Loss of food security (fisheries collapse, crop failures)
- Loss of clean water (aquifers depleted, rivers dry)
- Loss of clean air (pollution, smoke)
- Loss of stable climate (extreme weather, drought, flood)
- Loss of medicines (many drugs come from plants; extinctions destroy sources)
- Loss of meaning (the sacred, the wild, the other)

\section{Why Extraction Accelerates}

\subsection{Economic Model}

Modern economies measure growth as extraction. GDP counts the value of resources extracted, not the cost of depletion.

Cutting a forest generates revenue. The value of the standing forest (carbon storage, water filtration, species habitat) is not counted. So extraction is always profitable, replacement is never.

\subsection{Inequality}

The wealthy extract resources from poor countries and poor regions. They don't live with the consequences. Pollution happens downwind and downstream. Someone else's children drink contaminated water.

Cost is externalized. Profit is captured.

\subsection{Time Discounting}

A forest can be cut now (profit today) or preserved for 100 years of sustainable timber (profit spread over generations). Given interest rates and human psychology, cutting now is always rational economically, even if it's irrational ecologically.

\subsection{Commons Tragedy}

If a fishery is shared and unregulated, each fisher has incentive to extract as much as possible before others do. Tragedy: the fishery collapses because no one is disciplined enough to sustain it.

\section{The Return Cost}

To return to $S_0$ requires:

\subsection{Stop Extraction}

Halt fishing, logging, mining at current rates. Cost: loss of resources (timber, fish, metals) and loss of jobs.

\subsection{Restore}

Replant forests, restore fisheries, remove dams, restore wetlands, remove invasive species. Cost: capital, labor, time (decades to centuries).

\subsection{Reduce Consumption}

The primary driver of extraction is consumption. Return requires living with less material stuff, less energy, less convenience. Cost: reduction in material living standard (from wasteful to merely comfortable).

\subsection{Accept Consequences}

Stopping extraction now doesn't prevent all harms already in motion (climate change will worsen for decades). Ecosystems that are already collapsed can't be restored. Species that are extinct are gone. Cost: grief, loss, acceptance that some damage is permanent.

\section{Why Return Fails}

\subsection{Denial}

"The problem isn't that bad." "Technology will solve it." "It's too costly to change." Belief in denial buys time — time spent not returning, deepening displacement.

\subsection{Coordination Problem}

No single country can save the world. If one stops extraction and others don't, the saving country just loses economically while extraction continues elsewhere.

This requires global coordination. Nations won't coordinate if it costs them growth. Growth requires extraction. Return requires degrowth.

\subsection{Locked-In Systems}

Entire economies are built on extraction. Cars require oil. Electricity requires coal or dams. Agriculture requires synthetic fertilizers from fossil fuels. Food systems require monocultures and pesticides.

Changing this system is revolutionary. It requires redesigning civilization.

\section{The Likely Path}

Honest assessment: most of humanity will not return to sustainable extraction before the system collapses.

Instead:
1. Extraction continues until the resource is gone
2. Ecosystem collapses (fisheries, forests, aquifers, pollinators)
3. Collapse cascades (food, water, stability)
4. Societies fail (migration, conflict, famine)
5. Survivors rebuild on smaller scale

This is not certainty, but it is the probable path given current trajectories and lack of political will.

\section{Paths for Those Who Can}

Those who see the displacement and have resources can:

\subsection{Local Sustainability}

Grow food locally. Use renewable energy. Harvest water. Build community resilience. Live as if extraction will stop (because it will).

\subsection{Stewardship}

Restore land you have access to. Plant trees, rebuild soil, support local species. Be a guardian, not an extractor.

\subsection{Advocacy}

Vote, organize, protest for policies that slow extraction and support return. It's unlikely to work at scale, but it's required for dignity.

\subsection{Prepare for Descent}

Learn skills that will be valuable when growth stops: food production, basic medicine, tool repair, community building. Be useful in the world that's coming.

\section{Conclusion}

The displacement is from sustainable ecosystem to extractive collapse. The cost is borne by all life. The return cost is impossibly high. Most will not pay it.

Those who can should begin now. Reduce extraction, build resilience, support restoration, prepare for the transition.

The path back to $S_0$ is long. But it's the only path that leads anywhere worth going.

\end{document}
